\newcommand{\name}{{Sophia}}
\begin{abstract}
  Given the massive cost of language model pre-training,
  a non-trivial improvement of the optimization algorithm would lead to a material reduction on the time and cost of training.
  Adam and its variants have been state-of-the-art for years, and more sophisticated second-order (Hessian-based) optimizers often incur too much per-step overhead. In this paper, we propose \ours, \textbf{S}econd-\textbf{o}rder Cli\textbf{p}ped Stoc\textbf{h}astic Opt\textbf{i}miz\textbf{a}tion, a simple scalable second-order optimizer that uses a light-weight estimate of the diagonal Hessian as the pre-conditioner. The update is the moving average of the gradients divided by the moving average of the estimated Hessian, followed by element-wise clipping. The clipping controls the worst-case update size and tames the negative impact of non-convexity and rapid change of Hessian along the trajectory.
  Sophia only estimates the diagonal Hessian every handful of iterations, which has negligible average per-step time and memory overhead. On language modeling with GPT models of sizes ranging from 125M to 1.5B, Sophia achieves a 2x speed-up compared to Adam in the number of steps, total compute, and wall-clock time, achieving the same perplexity with 50$\%$ fewer steps, less total compute, and reduced wall-clock time. 
  Theoretically, we show that Sophia, in a much simplified setting, adapts to the heterogeneous curvatures in different parameter dimensions, and thus has a run-time bound that does not depend on the condition number of the loss. 
\end{abstract}


\section{Introduction}
Language models (LLMs) have gained phenomenal capabilities as their scale grows~\citep{radford2019language,kaplan2020scaling,brown2020language,  zhang2022opt,   
 touvron2023llama,openai2023gpt}. However, pre-training LLMs is incredibly time-consuming due to the massive datasets and model sizes---hundreds of thousands of updates to the model parameters are required. For example, PaLM was trained for two months on 6144 TPUs, which costed 10 million dollars~\citep{chowdhery2022palm}. 

Pre-training efficiency is thus a major bottleneck in scaling up LLMs. 
This work aims to improve pre-training efficiency with a faster optimizer, which either reduces the time and cost to achieve the same pre-training loss, or alternatively achieves better pre-training loss with the same budget.

Adam~\citep{kingma2014Adam} (or its variants~\citep{loshchilov2017decoupled,shazeer2018adafactor,you2019large}) is the dominantly used optimizer for training LLMs, such as GPT~\citep{radford2019language,brown2020language}, OPT~\citep{zhang2022opt}, Gopher~\citep{rae2021scaling} and LLAMA~\citep{touvron2023llama}. 
Designing faster optimizers for LLMs is challenging. First, the benefit of the first-order (gradient-based) pre-conditioner in Adam is not yet well understood~\citep{ liu2020understanding,zhang2020adaptive,kunstner2023noise}. Second, the choice of pre-conditioners is constrained because we can only afford light-weight options whose overhead can be offset by the speed-up in the number of iterations. For example, the block-diagonal Hessian pre-conditioner in K-FAC is  expensive for LLMs~\citep{martens2015optimizing,grosse2016kronecker, ba2017distributed,martens2018kronecker}. On the other hand, ~\citet{chen2023symbolic} automatically search among the light-weight gradient-based pre-conditioners and identify Lion, which is substantially faster than Adam on vision Transformers and diffusion models but only achieves limited speed-up on LLMs~\citep{chen2023symbolic}. 
 

This paper introduces \ours, \textbf{S}econd-\textbf{o}rder Cli\textbf{p}ped Stoc\textbf{h}ast\textbf{i}c Optimiz\textbf{a}tion, a light-weight second-order optimizer that uses an inexpensive stochastic estimate of the diagonal of the Hessian as a pre-conditioner and a clipping mechanism to control the worst-case update size. On pre-training language models such as GPT-2, {\ours} achieves the same validation pre-training loss with 50$\%$ fewer number of steps than Adam. Because {\ours} maintains almost the memory and average time per step, the speedup also translates to 50$\%$ less total compute and 50$\%$ less wall-clock time (See Figure~\ref{fig:head} (a)\&(b)). We also note that comparing the run-time to achieve the same loss is a correct way to compare the speed of optimizers for LLMs; see Section~\ref{sec:evaluation} for more details. 




Moreover, the scaling law based on model size from 125M to 770M is in favor of Sophia over Adam---the gap between Sophia and Adam with 100K steps increases as the model size increases (Figure~\ref{fig:head} (c)). 
In particular, Sophia on a 540M-parameter model with 100K steps gives the same validation loss as  Adam on a 770M-parameter model with 100K steps. Note that the latter model needs 40\% more training time and 40\% more inference cost.  

\begin{figure}[t]
\begin{center}
\includegraphics[width=0.4\textwidth]{figures/medium_2x_flops_plus_2309_arxiv.pdf}
\hspace{15pt}
\includegraphics[width=0.4\textwidth]{figures/large_2x_plusp_2309_all_arxiv.pdf}
\includegraphics[width=0.4\textwidth]{figures/1.5B_2x_200k_full_arxiv.pdf}
\hspace{15pt}
\includegraphics[width=0.4\textwidth]{figures/scale_plus_0923_arxiv.pdf}
\caption{{{\ours} achieves a 2x speedup over AdamW in GPT-2 pre-trained on OpenWebText and GPT NeoX pre-trained on the Pile. (a) (b) (c) Comparison of the number of steps needed to achieve the same level of validation loss on (a) GPT-2-medium (355M), (b) GPT-2-large (770M) and (c) GPT NeoX 1.5B. Across all model sizes, {\ours} needs 50$\%$ less time to reach the same validation loss as AdamW. (d) Validation losses of models with different sizes pre-trained for 100K steps. The gap between {\ours} and AdamW gets larger as models size grows. Notably, using {\ours} on a 540M-parameter model for 100K steps results in the same loss as using AdamW on a 770M-parameter model for 100K steps. See Section~\ref{sec:experiments} for details and more results. \label{fig:head}
}
}
\end{center}
\end{figure}

\begin{figure}
\begin{minipage}[b]{.46\textwidth}
\begin{minipage}[b]{\textwidth}
\begin{algorithm}[H]
	\caption{$\textup{Hutchinson}(\theta)$}
	\label{alg:sub}
	\begin{algorithmic}[1]
		\STATE {\bfseries Input:} parameter $\theta$.
            \STATE Compute mini-batch loss $L(\theta)$.
		\STATE Draw $u$ from $\mathcal{N}(0,\mathrm{I}_d)$.
        \RETURN $u \odot \nabla (\langle\nabla L(\theta), u\rangle)$.
	\end{algorithmic}
\end{algorithm}
\end{minipage}
\centering
\includegraphics[width=1.0\textwidth]{figures/toy_new.png}
\vspace{-20pt}
\caption{The motivating toy example. $\theta_{[1]}$ is the sharp dimension and $\theta_{[2]}$ is the flat dimension. GD's learning rate is limited by the sharpness in $\theta_1$, and makes slow progress along $\theta_{[2]}$. Adam and SignGD bounce along $\theta_{[1]}$ while making slow progress along $\theta_{[2]}$. Vanilla Newton's method converges to a saddle point. {\ours} makes fast progress in both dimensions and converges to the minimum with a few steps.\label{fig:toy}}

\vspace{-10pt}

\end{minipage}\hfill
\begin{minipage}[b]{.49\textwidth}
\begin{minipage}[b]{\textwidth}
\begin{algorithm}[H]
	\caption{ $\textup{Gauss-Newton-Bartlett}(\theta)$}
	\label{alg:sub_gn}
	\begin{algorithmic}[1]
		\STATE {\bfseries Input:} parameter $\theta$.
		\STATE Draw a mini-batch of input $\{x_b\}_{b=1}^{B}$.
            \STATE Compute logits on the mini-batch: $\{f(\theta, x_b)\}_{b=1}^{B}$.
            \STATE Sample $\hat{y}_b\sim \textup{softmax}(f(\theta, x_b)), \forall b\in [B]$.
            \STATE Calculate $\hat{g} = \nabla(1/B\sum\ell(f(\theta, x_b),\hat{y}_b))$.
        \RETURN $B\cdot \hat{g} \odot \hat{g}$.
	\end{algorithmic}
\end{algorithm}
\end{minipage}
\begin{minipage}[b]{\textwidth}
\begin{algorithm}[H]
	\caption{\ours~}
	\label{alg:hess_opt}
	\begin{algorithmic}[1]
		\STATE {\bf Input:} $\theta_1$, learning rate $\{\eta_t\}_{t=1}^T$,  hyperparameters $\lambda,\gamma, \beta_{1},\beta_{2}$, $\epsilon$, and estimator choice Estimator $\in\{$Hutchinson, Gauss-Newton-Bartlett$\}$\STATE Set $m_{0} = 0$, $v_{0} = 0$, $h_{1-k} = 0$
		\FOR{$t=1$ {\bf to} $T$}
		\STATE Compute minibach loss $L_t(\theta_t)$.
        \STATE Compute $g_t = \nabla L_t(\theta_t)$.
		\STATE  $m_{t} = \beta_{1} m_{t-1} + (1 - \beta_{1}) g_{t}$ 
        \IF{ $t\ \mathrm{mod}\  k = 1$}
        \STATE Compute $\hat{h}_{t} = \textup{Estimator}(\theta_t)$.
		\STATE  $h_t = \beta_2 h_{t-k} + (1 - \beta_2) \hat{h}_{t}$
        \ELSE \STATE  $h_t = h_{t-1}$
        \ENDIF
        \STATE $\theta_t = \theta_t - \eta_t\lambda \theta_t$ (weight decay)
		\STATE $\theta_{t+1} = \theta_{t} - \eta_t \cdot \clip(m_t / \max\{\gamma \cdot h_t,\epsilon\},1)$
		\ENDFOR
	\end{algorithmic}
\end{algorithm}
\end{minipage}
\end{minipage}
\end{figure}


Concretely, {\ours} estimates the diagonal entries of the Hessian of the loss using a mini-batch of examples every $k$ step (with $k = 10$ in our implementation). 
We consider two options for diagonal Hessian estimators: (a) an unbiased estimator that uses a Hessian-vector product with the same run-time as a mini-batch gradient up to a constant factor, and (b) a biased estimator that uses one mini-batch gradient calculated with resampled labels. Both the two estimators only introduce 5\% overheads per step (in average).
At every step, {\ours} updates the parameter with an exponential moving average (EMA) of the gradient divided by the EMA of the diagonal Hessian estimate, subsequently clipped by a scalar. (All operations are element-wise.) See Algorithm~\ref{alg:hess_opt} for the pseudo-code.

Additionally, {\ours} can be seamlessly integrated into existing training pipelines, without any special requirements on the model architecture or computing infrastructure. With the either of the Hessian estimators, {\ours} only require either standard mini-batch gradients, or Hessian-vector products which are supported in auto-differentiation frameworks such as PyTorch~\citep{pytorch2019} and JAX~\citep{jax2018github}. 


Thanks to the Hessian-based pre-conditioner, {\ours} adapts more efficiently, than Adam does, to the heterogeneous curvatures in different parameter dimensions, which can often occur in the landscape of LLMs losses and cause instability or slowdown. 
Sophia has a more aggressive pre-conditioner than Adam---{\ours} applies a stronger penalization to updates in sharp dimensions (where the Hessian is large) than the flat dimensions (where the Hessian is small), ensuring a uniform \textit{loss} decrease across all parameter dimensions. In contrast, Adam's updates are mostly uniform across all parameter dimensions, leading to a slower loss decrease in flat dimensions. (See Section~\ref{sec:motivation} for more discussions.) 
These make Sophia converge in fewer iterations. Thanks to the light-weight diagonal Hessian estimate, the speed-up in the number of steps translates to a speed-up in total compute and wall-clock time. 



\ours's clipping mechanism controls the worst-case size of the updates in all directions, safeguarding against the negative impact of inaccurate Hessian estimates, rapid Hessian changes over time, and non-convex landscape (with which the vanilla Newton’s method may converge to local maxima or saddle points instead of local minima). The safeguard allows us to estimate Hessian infrequently (every $k=10$ step) and stochastically. In contrast, prior second-order methods often update Hessian estimates every step~\citep{martens2015optimizing,grosse2016kronecker,  anil2020scalable,yao2021adahessian}. 


We provide theoretical analyses of much simplified versions of Sophia on convex functions. 
The runtime bound does not depend on the local condition number (the ratio between maximum and minimum curvature at the local minimum) and the worst-case curvature (that is, the smoothness parameter), demonstrating the advantage of Sophia in adapting to heterogeneous curvatures across parameter dimensions. 